\subsection{Supplied Choices}
The contract states what may change and what must be preserved. Its inputs are component models and transfers, requirements and their initializers, and fixed endpoint controllers. Synthesis supplies the coordination among them.

\paragraph{Finite transition systems and safety monitors.}
A labeled transition system (LTS) $\mathcal L=(S,\Sigma_{\mathcal L},\delta,s^0)$ has nondeterministic transitions $\delta:S\times\Sigma_{\mathcal L}\to2^S$. Partition events as $\Sigma=\Sigmauc\uplus\Sigmac$: controllers may disable controllables but cannot select outcomes or suppress uncontrollable (UC) events. A synchronous-product event requires a transition in every component whose alphabet contains it; other components retain their state~\cite{MageeKramer2006}.
A monitor is an automaton that checks a requirement as events occur. Requirement $r$ has a total deterministic monitor $t_r=(M_r,\Sigma,\eta_r,m_r^0,F_r)$ with absorbing errors for violating prefixes~\cite{AlpernSchneider1985,KupfermanVardi2001}. The empty monitor product has one non-error state. Monitors observe labels; policies also observe successor component states.

\label{def:fg-inputs}
\label{def:update_plan}
\paragraph{Components and transfers.}
For each $i\in I$, let $e_i^{\mathrm{old}},e_i^{\mathrm{new}}$ be the pre-update and post-update components and let $g_i\subseteq S_{e_i^{\mathrm{old}}}\times S_{e_i^{\mathrm{new}}}$ be their state-transfer relation.
  Define $g_i(s)=\{s'\mid(s,s')\in g_i\}$. The supplied $g_i$ includes every modeled transformation outcome. Let $E^v=\mathop{\parallel}_{i\in I}e_i^v$.

\paragraph{Requirements and activation.}
Let pairwise-disjoint $R^{\mathrm{old}},R^{\mathrm{new}},R^{\mathrm{upd}}$ contain the pre-update, post-update, and update-interval safety requirements, respectively.
For each $r\in R^{\mathrm{new}}\cup R^{\mathrm{upd}}$, the input supplies a partial \emph{monitor-initialization map} $\iota_r$ from version-tagged component tuples to $M_r\setminus F_r$. Its domain specifies where activation is allowed.
Section~\ref{sec:monitor-meaning} defines history-based and activation-scoped interpretations. For $r\in R^{\mathrm{upd}}$, $\operatorname{dom}\iota_r$ includes every reachable old component state.

\paragraph{Endpoints and handover targets.}
\label{def:endpoint_controllers}
For $v\in\{\mathrm{old},\mathrm{new}\}$, the supplied finite controller $C^v$ preserves uncontrollables and makes $E^v\parallel C^v$ deadlock-free and satisfy $R^v$. Hiding controller auxiliary events leaves observation alphabet $\bigcup_i\Sigma_{e_i^v}$; the union over both versions is the ordinary alphabet $\SigmaN$. Write $\Reach$ for initial-state reachability and
\[
\mathsf{CL}^v=E^v\parallel C^v\parallel\mathop{\parallel}_{r\in R^v}t_r.
\]
Monitors are passive: every endpoint transition originates in $E^v\parallel C^v$; a monitor only updates on that label and adds no enabled event.
The input supplies nonempty, checked targets $Z_{\mathrm{load}}\subseteq\Reach(\CLnew)$ with component, controller and new-monitor states.

\paragraph{Update commands and order.}
The set of update events, each executed exactly once, is
  \begin{equation}
  \SigmaU=
    \{\rho_i\mid i\in I\}
    \cup\{\ustop{r}\mid r\in R^{\mathrm{old}}\}
    \cup\{\ustart{r}\mid r\in R^{\mathrm{new}}\}.
  \end{equation}
These controllable events transfer components, stop old monitors or start new ones, respectively, and are disjoint from ordinary events $\SigmaN$:
$\Sigma=\SigmaN\uplus\SigmaU$ and $\SigmaU\subseteq\Sigmac$. Each transfer event $\rho_i$ is one atomic model transition: the controller chooses the update, but cannot choose its $g_i$ outcome.

An optional strict partial order $\prec\ \subseteq\SigmaU\times\SigmaU$ expresses precedence among update operations. For a pending-event set $P$, event $a$ is order-enabled exactly when $a\in P$ and every $a'\prec a$ has left $P$.
Monitors self-loop on update events that they do not specify.
\paragraph{Eligible contracts.}
An eligible contract satisfies the finiteness, typing, totality and non-error conditions above; dependencies form a strict partial order, and endpoint controller alphabets cover environment labels. Eligibility is structural; Section~\ref{sec:monitor-meaning} justifies initializer meanings.

\Needspace{4\baselineskip}
\subsection{States, Transitions, and Completion}
\label{sec:game-rules}

\label{def:configuration}
\label{def:init}
\emph{States and all entry states.}
An update state $q=(\xi,\mu,P)$ records each component's version/state $\xi(i)=(v,s)$ with $s\in S_{e_i^v}$, each monitor's state $\mu(t_r)\in M_r\cup\{\bot\}$ ($\bot$ means inactive), and pending commands $P\subseteq\SigmaU$.
Here $\rho_i\in P$ exactly when component $i$ is old. Old monitors remain active until their stop leaves $P$; new monitors become active when their start leaves $P$; update monitors are always active. 

Admission $\Init(z)$ copies $z$'s component and old-monitor states, installs interval initializers, leaves new monitors inactive and sets $P=\SigmaU$. The update policy replaces $C^{\mathrm{old}}$ and its scheduling restrictions. Declared old monitors remain until their stops; undeclared objectives, including old liveness, are not preserved. Duties needed throughout belong in interval monitors. The entry set is
\begin{equation}
 Q_0=\{\Init(z)\mid
 z\in\Reach(\mathsf{CL}^{\mathrm{old}})\}.
\label{eq:initial}
\end{equation}
All-entry coverage forbids waiting for a favorable old state.
\par\noindent\emph{Four kinds of induced transitions.}
\label{def:successor}
For $\xi(i)=(v,s)$, component $i$ uses $e_i^v$. The active-component product determines enabled ordinary uncontrollables $U(q)$, including those with unsafe outcomes, independently of updates and goals. For every $a\in\SigmaU$, impose
\begin{equation}
 U(q)\ne\emptyset\quad\Longrightarrow\quad\Post(q,a)=\emptyset.
 \label{eq:nonpreemption}
\end{equation}
Thus even a local command waits until the active-component product has no enabled uncontrollable event. A finite operation may drain to such a point; perpetual polling may prevent it. This global \emph{uncontrollable quiescence} differs from transaction quiescence~\cite{KramerMagee1990Quiescence} and tranquility~\cite{VandewoudeEtAl2007Tranquility}. Ordinary controls may coexist with UC. The following rules define $\Post$; unmentioned coordinates stay unchanged. ``Advance'' means $\mu'(t_r)=\eta_r(\mu(t_r),a)$ for the specified active monitors.

\begin{description}
  \item[Ordinary event $a\in\SigmaN$.]
  If $a$ belongs to the local event set of at least one active component and every active component having $a$ can transition, take every synchronous-product outcome and advance every active monitor.

  \item[Component update $\rho_i$.]
  If $U(q)=\emptyset$, order-enabled, and $\xi(i)=(\mathrm{old},s)$, set $\xi'(i)=(\mathrm{new},s')$ for each $s'\in g_i(s)$, delete $\rho_i$ from $P$, and advance every active monitor.

  \item[Requirement deactivation $\ustop{r}$.]
  If $U(q)=\emptyset$ and order-enabled, advance every active monitor except $t_r$, retain inactive monitors at $\bot$, set $\mu'(t_r)=\bot$, and delete $\ustop{r}$ from $P$; $t_r$ does not observe it.

  \item[Requirement activation $\ustart{r}$.]
  If $U(q)=\emptyset$, order-enabled, and $\xi\in\operatorname{dom}\iota_r$, advance every active monitor, retain other inactive monitors at $\bot$, install $\mu'(t_r)=\iota_r(\xi)$, and delete $\ustart{r}$ from $P$; $t_r$ begins online observation afterward.
\end{description}
If a precondition is false, the successor set is empty; every nondeterministic outcome is included and environment-chosen.
Let $\mathcal Q$ contain every state reachable from $Q_0$ through these transitions. 
\par\noindent\emph{Safety and handover-ready goals.}
\label{def:safe_goal}
An update state $q=(\xi,\mu,P)$ is safe, written $\Safe(q)$, exactly when every active monitor is outside its error states.
For $P=\emptyset$ and $z\in Z_{\mathrm{load}}$, write $q\sim z$ when component states (without version tags) and all new-monitor states agree. The handover map $\kappa$ selects such a full endpoint state $z$. Goals are
\begin{equation}
\Goal=\{q=(\xi,\mu,\emptyset)\in\mathcal Q\mid
\Safe(q)\land U(q)=\emptyset\land\exists z\in Z_{\mathrm{load}}:q\sim z\}.
\label{eq:goal}
\end{equation}
The goal is a quiescent handover boundary, not merely completion of the command list. At a goal, interval obligations end and matched new monitors continue. Handover disables ordinary controllables and installs $\kappa(q)$'s controller state, outside $\Sigma$ and the rank. UC cannot compete; compatible targets may have different continuations.
\par\noindent\emph{State-dependent policies and finite completion.}
\label{def:strategy}
A memoryless update policy is a function
$\sigma:\mathcal Q\setminus\Goal\to2^{\Sigmac}$ selecting only events with nonempty successor sets.

The possible next states under this policy are
\begin{equation}
 \SuccSet_\sigma(q)=
 \bigcup_{u\in U(q)}\Post(q,u)
 \ \cup\
 \bigcup_{a\in\sigma(q)}\Post(q,a).
\label{eq:strategy-successors}
\end{equation}
The environment chooses any uncontrollable and any permitted-event outcome. The returned policy retains one controllable at a quiescent state.

Let $\Reach_\sigma(Q_0)$ be the policy closure from $Q_0$, stopping at goals. 

\emph{Runs, winning, and witnesses.}
A run follows the policy successors and stops at a goal. It is \emph{maximal} if infinite or ending at a goal or a state without successors. A single policy $\sigma$ is winning for an FG-DUCS instance if, from every $q_0\in Q_0$,
(i) every reachable update state is safe, (ii) every reachable non-goal state has a successor, and
(iii) every maximal update run reaches $\Goal$ after finitely many events.
This \emph{finite completion} counts events and assumes no fairness. Winning policies exclude controllable self-loops during update; a reached UC self-loop makes a state losing. Ordinary actions can progress, as Cell's moves and inspection do. Handover restores the supplied controller's full matched behavior.

A \emph{progress-ranking function} for a winning policy is $\rank:\Reach_\sigma(Q_0)\to\mathbb N$, with value 0 at goal states and satisfying $\rank(q')<\rank(q)$ for every non-goal state $q$ and every
$q'\in \SuccSet_\sigma(q)$.

\paragraph{The FG-DUCS problem.}
Decide whether one policy safely completes the eligible game from every entry under every permitted event and outcome. Return $(\sigma,\rank,\kappa)$ covering all $Q_0$, or a losing-region certificate.
